% !TEX root = updates.tex
\begin{addblock}
\subsection{Hidden escalation: each side's belief over the level}\label{sec:hidden}
In the game so far both sides see $\ell^t$. Here we drop that assumption. Each side
sees incidents, not the level. As an option (\texttt{--hidden}; off by default, which reproduces every other
result exactly) we remove $\ell^t$ from both players' observations. The engine still runs the true level with the same
hazard \eqref{eq:hazard}, escalation rule and payoffs \eqref{eq:stage}--\eqref{eq:payoff}. Each player instead sees the
zones where incidents happened, $I^t\subseteq\C$ (public, $N^t=|I^t|$), and whether war has broken out. Its stage
payoffs are not shown to it during play. Against a fixed Red policy, Blue then faces a partially observable Markov
decision process (POMDP) in $\ell$ as well as in Red's positions. With both sides adapting, it is a partially
observable stochastic game.

\textbf{Belief:} player $i$ holds $b^t_i(e)=\Pr(\ell^t=e\mid\text{$i$'s history})$, $e=0,\dots,L$, starting as a point
mass at $\ell^1$. Each period it is updated in the same order as the game:
\begin{align*}
&\text{(1) evidence:} && \tilde b(e)\propto b^t_i(e)\,\psi^t_i(e),\\
&\text{(2) transition:} && b'(e')=\textstyle\sum_e\tilde b(e)\,Q(e'\mid e,N^t),\qquad
Q(e'\mid e,N)=\begin{cases}\1[e'=\min(L,e+N)] & N\ge1,\\ \mu(e)\,\1[e'=e-1]+(1-\mu(e))\,\1[e'=e] & N=0,\end{cases}\\
&\text{(3) no war:} && b^{t+1}_i(e)\propto b'(e)\,\1[e<L]\quad\text{(a point mass at $L$ if war broke out)},
\end{align*}
with $\mu(e)$ as in Table~\ref{tab:params} ($\mu(0)=0$). Filter \textbf{F0} sets $\psi\equiv1$. It uses only the
incident count and ``no war yet'', and both sides hold the same belief. Filter \textbf{F1} also uses where incidents did
and did not happen, since a higher level makes every incident likelier. Over the zones $\C_i$ where $i$ has ships,
\[
\psi^t_i(e)=\prod_{v\in\C_i}\hat\rho_v(e)^{\1[v\in I^t]}\big(1-\hat\rho_v(e)\big)^{\1[v\notin I^t]},
\]
where $\hat\rho_v(e)$ is the incident probability of \eqref{eq:hazard} at level $e$, computed from $i$'s own ships and
postures and its believed opponent counts and postures. Under F1 the two sides' beliefs differ. The filters draw no
random numbers, so games stay on common random numbers. We checked both against brute-force enumeration of every
path of the level (agreement to $10^{-15}$).

\textbf{Deciding with a belief:} wherever a policy used $\ell^t$ it now averages over the levels. Each level is
evaluated separately, never at the mean level, because the hazard is not linear in $\ell$ (D3 also switches at $\ell=3$):
\[
\bar\rho_v=\textstyle\sum_e b(e)\big(1-e^{-h_v(e)}\big),\qquad
\bar\varepsilon_i=\sum_e b(e)\,\varepsilon_i\,(1-e/L).
\]
The MIP of \S\ref{sec:mip} keeps its form and size. $\bar\rho_v$ replaces $\rho_v$ in the option values and in the
risk budget, whose right side becomes $\bar\varepsilon_i+\xi$. A belief is a point mass at the true level when $\ell$
is seen, which gives back the original MIP. In the rule-based and template policies a test $\ell\ge k$ becomes
$\Pr_b(\ell\ge k)\ge\tau$, with threshold $\tau=0.25$. We also log each plan's believed chance of war this period,
$\sum_{e<L}b(e)\Pr(N\ge L-e\mid e)$, with $N$ the number of incidents (independent zones at the believed opponent
positions). An optional cap on it makes MIP-BR solve again with half the budget, up to four times.
\textbf{What this is not:} policies still look one period ahead on the current belief. No one plans to learn about
the level, and Red does not model Blue's belief.
\end{addblock}
